\chapter{Alt-screen and terminal state}
\label{ch:altscreen}

\begin{aside}
The terminal is shared global state. Your TUI is one tenant
among many. Leaving it in raw mode, with mouse on, on the
alternate screen, is a sin against the next user.
\end{aside}

\section{Two screens}

A modern terminal has two text buffers:

\begin{itemize}[leftmargin=*]
  \item The \textbf{main} screen, where the shell prompt and
    command output live.
  \item The \textbf{alternate} screen, a separate scrollback-
    free buffer for full-screen apps.
\end{itemize}

Vim, htop, less, tmux --- they all switch to the alternate
screen, do their work, and switch back. The user's shell
history is preserved.

\maya{} enters the alternate screen on startup
(\code{ESC [ ? 1049 h}) and exits on shutdown
(\code{ESC [ ? 1049 l}).

\section{Why 1049 and not 47}

There are two ``alternate screen'' modes:

\begin{itemize}[leftmargin=*]
  \item \code{? 47}: just switch buffers.
  \item \code{? 1049}: save cursor, switch buffers, clear the
    alternate, restore cursor on exit.
\end{itemize}

\code{1049} is the modern choice and handles the corner cases
(cursor position preserved across the switch). \maya{} uses it.

\section{Terminal modes}

A terminal has many modes. Each is set by an escape sequence;
together they form the ``terminal state''. The relevant ones
for a TUI:

\begin{itemize}[leftmargin=*]
  \item \textbf{Raw mode (tcsetattr).} No line buffering, no
    echo, no signal generation on Ctrl-C. Set via the
    \code{termios} struct.
  \item \textbf{Alt screen.} \code{? 1049}.
  \item \textbf{Mouse mode.} \code{? 1000}, \code{? 1002},
    \code{? 1003}, \code{? 1006}.
  \item \textbf{Bracketed paste.} \code{? 2004}.
  \item \textbf{Cursor visibility.} \code{? 25}.
  \item \textbf{CSI-u.} \code{> 1 u}.
  \item \textbf{Focus events.} \code{? 1004}.
\end{itemize}

\section{The boot bundle}

\maya{} emits a fixed bundle on startup:

\begin{lstlisting}
ESC [ ? 1049 h       enter alt screen
ESC [ ? 25   l       hide cursor
ESC [ ? 1006 h       SGR mouse
ESC [ ? 1002 h       mouse press+drag
ESC [ ? 2004 h       bracketed paste
ESC [ ? 1004 h       focus events
ESC [ > 1 u          CSI-u
ESC [ 2 J            clear screen
ESC [ H              cursor home
\end{lstlisting}

Order matters: clear after alt-screen, otherwise you cleared
the main screen. Cursor home so the first frame's positioning
is predictable.

\section{The shutdown bundle}

The reverse, in reverse order:

\begin{lstlisting}
ESC [ < u            CSI-u off
ESC [ ? 1004 l       focus events off
ESC [ ? 2004 l       bracketed paste off
ESC [ ? 1002 l       mouse modes off
ESC [ ? 1006 l
ESC [ ? 1000 l
ESC [ ? 25 h         cursor visible
ESC [ 0 m            reset SGR
ESC [ ? 1049 l       exit alt screen
\end{lstlisting}

The SGR reset is important: leaving in red-on-yellow would
colour the user's shell.

\section{Crash safety}

If the program terminates without running shutdown (signal,
\code{std::terminate}, hard kill), the terminal is left in
whatever state. \maya{} mitigates by:

\begin{itemize}[leftmargin=*]
  \item Installing a \code{std::atexit} hook that emits the
    shutdown bundle.
  \item Installing signal handlers for SIGTERM, SIGINT,
    SIGHUP that emit shutdown and re-raise.
  \item RAII guards on the \code{Terminal} object that emit
    shutdown in the destructor.
\end{itemize}

A \code{kill -9} bypasses all of these. The user can recover
with \code{reset} or by closing the terminal window.

\section{The reset command}

On Unix, \code{reset} (or \code{stty sane}) resets the
terminal to a known-good state: cooked mode, echo on, alt
screen off, all mouse modes off, all colours reset, all
private modes off.

It is the universal ``unstuck my terminal'' tool. Add a
mention in your app's docs.

\begin{tryit}
Run a \maya{} app, suspend with Ctrl-Z, and try to use the
shell. Mouse clicks fire as escape sequences. Resume with
\code{fg}, then exit normally. Compare to crashing the app
mid-run; that needs a \code{reset}.
\end{tryit}

\section{tmux and screen}

Inside tmux or GNU screen, the terminal you talk to is the
multiplexer, not the outer terminal. The multiplexer
forwards the modes selectively. Some modes (CSI-u, focus
events) are passed through only if you opt in.

\maya{}'s requests are advisory. The multiplexer chooses.

\section{Terminal queries}

You can ask the terminal about its state:

\begin{itemize}[leftmargin=*]
  \item \code{ESC [ 6 n}: report cursor position.
  \item \code{ESC [ ? 25 6 n}: report cursor position
    extended (some terminals).
  \item \code{ESC [ 14 t}: report window size in pixels.
  \item \code{ESC [ 18 t}: report window size in cells.
  \item \code{ESC [ > q}: report terminal name.
\end{itemize}

The replies arrive on stdin like any other input. \maya{}'s
parser recognises them and routes them to the query handler
instead of the key dispatch.

Queries are useful for cell-pixel size (for inline image
sizing) and for detecting the terminal at startup. Use
sparingly: a query that goes unanswered (terminal does not
support it) is just dead bytes.

\section{The Da1 query}

\code{ESC [ c} asks ``which terminal are you?''. The reply is
a sequence like \code{ESC [ ? 1 ; 2 c} where the numbers are
DEC feature flags.

\maya{} sends Da1 at startup as a smoke test: if no reply
arrives within 200 ms, the terminal is suspect (a pipe to a
file, perhaps, not a real terminal). The library falls back
to non-interactive mode.

\section{Recap}

\begin{recap}
\begin{itemize}[leftmargin=*]
  \item Two screens; \maya{} uses the alternate.
  \item Raw mode plus a dozen private modes form the boot
    bundle.
  \item Shutdown is reverse order; SGR reset matters.
  \item Crash safety via atexit, signals, RAII.
  \item \code{reset} unsticks anything.
  \item Queries get replies on stdin; the parser routes
    them.
\end{itemize}
\end{recap}

\section*{Workshop}

\begin{enumerate}[leftmargin=*]
  \item Print the Da1 reply for your terminal. Compare to
    xterm, kitty, alacritty, Apple Terminal.
  \item Crash your app deliberately (assert, segfault).
    Confirm the terminal is left usable thanks to the
    atexit hook.
  \item Run your app inside tmux. Try CSI-u and confirm it
    works only after configuring tmux's
    \code{extended-keys} option.
\end{enumerate}
